\begin{table*}[t]
\caption{Comparison of Environmental Inequality and Spatial Air Quality Mapping Methodologies}
\label{tab:literature_comparison}
\centering
\footnotesize
\resizebox{\textwidth}{!}{%
\begin{tabular}{p{3.5cm}p{2.8cm}p{2.8cm}p{2.8cm}p{4.2cm}}
\toprule
\textbf{Methodology / Paradigm} & \textbf{Target Observational Data} & \textbf{Spatial Validation Scheme} & \textbf{Uncertainty Estimation} & \textbf{Primary Operational Characteristics} \\
\midrule
Econometric Surface Mapping~\cite{colmer2020,jbaily2022} & Satellite-derived PM$_{2.5}$ & Census-tract aggregation & Clustered standard errors & Large spatial coverage; relies on model-derived aerosol surfaces \\
Epidemiological Reanalysis Surfaces~\cite{di2017} & CTM / AOD reanalysis grids & Regional grid comparison & Ensemble spread & Continuous fields; coarse resolution smooths localized monitor-scale gradients \\
Classical Geostatistics (IDW/Kriging)~\cite{shepard1968,cressie2011} & Physical in-situ network & Distance-based interpolation & Kriging variance / empirical & Strong local distance-decay prior; lacks non-linear socioeconomic feature fusion \\
Tabular Gradient Boosting~\cite{meyer2018,roberts2017} & Station coordinates + census & Random cross-validation & Often unquantified & Captures feature interactions; risks spatial discontinuity across unmonitored clusters \\
\midrule
\textbf{Spatial Blocked Benchmark (This Study)} & \textbf{EPA AQS Ground Monitors} & \textbf{8-Cluster Spatially Blocked CV} & \textbf{Quantile Prediction Bounds} & \textbf{Direct in-situ grounding; evaluates domain shift to unmonitored regions} \\
\bottomrule
\end{tabular}%
}
\end{table*}
